\section{Antisocial Behavior}

One of Tomberg's guiding spirits in the French Hermetic Tradition was \textbf{Josephin Peladan}. The vision of the mass man, the unitiated man, is restricted to two dimensions. He absorbs his ideas and opinions from those around him and is subject to the vagaries of popular opinion. It can be said that he not so much lives his life but rather life lives through him. The Mage, or initiate, on the other hand, sees the third dimension of life, which is depth. Peladan ties Baptism and the grace of faith to this depth. Baptism that frees us from the hypnosis of Society. He writes, in How to Become a Mage:

\begin{quotex}
Baptism as an initiation rite makes us children of God, but Society dooms us to evil through is laws and education. Faith enlightens us, but it is in perpetual conflict with Society. The initiate, in order to make the grace of baptism full and effective, must renounce anew Society, its boundaries, its crimes, so that the fear of God makes him prefer the intimate greatness to the degrading favors of the country.
\end{quotex}


The uninitiated man takes the world as it appears to him as the norm; this is not the real world, but is a second reality, or a shadow world hiding its true source; such a man lives in that second reality. The initiate begins to wake up and must see through that second reality in order to become fully conscious. Of course, in our day and age, everyone thinks he is a rebel and boasts about flaunting societal mores. Unable to conceive of any higher goal, this faux rebel can only act out against the taboos against sex, drugs, and other anti-social behavior. This is just another trap, another way of being moulded. This type is similar to the Aghori of India\footnote{\url{https://gornahoor.net/?p=3161}}. Peladan explains:

\begin{quotex}
Before you think and choose, society takes over your being and moulds it, as is its right. Once you think and choose, remove those received imprints, that is to say, liberate yourself from contemporary habits, as is your right.
\end{quotex}


How many readers of Meditations on the Tarot are truly willing to remove their received imprints from society? Has anyone really changed his opinions on religion, politics, society, morality following a reading of Meditations on the Tarot? If so, it is quite rare, and usually not in the direction that \textbf{Valentin Tomberg} is pointing.

Peladan has a message, in particular, to Roman Catholics, in his discussion of the esoteric meaning of marriage. In particular, an initiate needs to integrate his tradition with Catholicism. In Peladan's own words:

\begin{quotationx}
The correspondence between marriage and magic, is the union of the initiate with the tradition contained in the chefs-d'oeuvre of Roman Catholicism, and the care of combining all the scattered morsels of truth and those belonging to religions that have disappeared.

The virtue of the initiate is formed from equanimity; the beatitude given to the peacemakers applies to him.

The highest work of mercy consists in making, in one's thought, a sepulcher to sublime thoughts; gathering the beautiful ideas lost in ancient books and, I say it especially to the Roman congregations, lifting up into their understanding a cenotaphe to Plato, rethinking the sublime thoughts of Confucius or of Zarathustra, will always be the highest of pieties as well as the rarest.
\end{quotationx}


So, now that the reader is ready to think and choose, what are the choices that Peladan offers us? They are:

\begin{quotex}
In order to choose, know that you have three destinies:

\begin{enumerate}
\item you can be an animal like that decadent superficial man called savage;

\item a soulish man, like everyone else;

\item a spiritual person like \textbf{St Thomas} or \textbf{Dante}.
\end{enumerate}

\begin{itemize}


\item Animal: be beautiful

\item soulish man: be good

\item spiritual man: seek the Grail.
\end{itemize}
\end{quotex}


This is the path from a life centered on the body to one centered on the soul to one centered in Spirit. Peladan gives us the magical formula to progress:

\begin{itemize}


\item In order to improve, soulify your instincts

\item in order to make yourself meek, spiritualize your feelings

\item in order to reach the absolute, develop within yourself abstraction.
\end{itemize}

Originally published on 22 October 2011 on medtarot.


\flrightit{Posted on 2023-03-07 by Cologero}

